% !TEX root = ../main.tex
% ESEMPIO 0c: slide dell'approccio (Overview, Details, diagramma).
%   - Overview: una frase sull'idea e un punto per componente; le slide dopo corrispondono ai punti.
%   - Details: un blocco con una tupla e la descrizione dei suoi elementi (\begin{block}); per un blocco
%     lungo si può rimpicciolire il testo con { \footnotesize ... } attorno al blocco.
%   - Diagramma: una figura da assets/tex/ dentro adjustbox, con poche parole.
%   Da copiare: la struttura del frame che serve, poi sostituire i CHANGE.
%   L'esempio (running example) si può raccontare dentro ogni componente, mentre lo presenti, oppure passo passo
%   dopo l'approccio. Dipende dal tempo.

\begin{frame}
    \frametitle{Example: approach overview}
    % GUIDA: una frase sull'approccio e sotto un punto per componente. Le slide dopo corrispondono ai punti, nello stesso ordine:
    % è un indice implicito dei contenuti, non una scaletta.

    Main \alert{idea}: CHANGE: one sentence stating the idea of the approach

    \vspace{1em}

    \begin{itemize}\setlength{\itemsep}{1em}
        \item CHANGE: aspect 1 of the idea, with the \alert{key words} highlighted
        \item CHANGE: aspect 2 of the idea
        \item CHANGE: aspect 3 of the idea
    \end{itemize}

\end{frame}

\begin{frame}
    \frametitle{Example: approach details}
    % GUIDA: una slide per componente, dall'alto verso il basso: dove sta nell'approccio complessivo, poi cos'è. Meglio un'immagine,
    % anche animata, di una definizione lunga. Si copia il frame per ogni componente (Details (1), (2), ...).
    \begin{block}{CHANGE: element name}
        A \alert{CHANGE: element} is a tuple \alert{$\angled{x_1, \{y_1,\ldots, y_n\}, \referred{X}, \signature}$}, where

        \vspace{1em}

        \begin{itemize}\setlength{\itemsep}{.5em}
            \item \alert{$x_1$} is CHANGE: description of element 1
            \item \alert{$\{y_i\}$} is a set of \emph{CHANGE: elements}, CHANGE: description of element 2
            \item \alert{$\referred{X}$$=$$\{\referred{x}_1,$ \ldots$,$ $\referred{x}_n\}$} is a set of \alert{\emph{referred elements}}, \alert{each} $\referred{x}_i$ \alert{CHANGE: relation to the element} (CHANGE: description of element 3)
            \item \alert{$\signature$} is CHANGE: description of element 4.
        \end{itemize}
    \end{block}

\end{frame}

\begin{frame}
    \frametitle{Example: state diagram}
    % GUIDA: un altro aspetto del processo, mostrato con un diagramma e poche parole (qui un diagramma di stati).

    \begin{adjustbox}{max totalsize={\linewidth}{\textheight},center}
        		\begin{tikzpicture}
			\tikzset{
				every node/.style={font=\scriptsize}
			}

			% \node[rboxes] (valid) {Valid};
			% \node[rboxes] (suspended) [right=2 of valid] {Suspended};
			% \node[rboxes] (revoked) [right=2 of suspended] {Revoked};
			\node[rboxes] (valid) {\certStatusValid{\cert}};
			\node[rboxes] (suspended) [right=2 of valid] {\certStatusSuspended{\cert}};
			\node[rboxes] (revoked) [right=2 of suspended] {\certStatusRevoked{\cert}};

			\draw[-latex] (valid) to [out=30,in=180-30] node [above]{\makecell{AL failure\\\referred{\cert} suspended}} (suspended);
			\draw[-latex] (suspended) to [out=180+30,in=-30] node [below]{\makecell{Issue remediated}} (valid);
			\draw[-latex] (valid.north) -- ++(0,1.5) -- node[midway,above]{\referred{\cert} revoked} ([yshift=1.5cm]revoked.north) -- (revoked.north);
			\draw[-latex] (suspended) -- node [above]{\makecell{Issue not\\remediated}} (revoked);
		\end{tikzpicture}
    \end{adjustbox}

\end{frame}
